\section{Conclusiones}
\label{sec:conclusiones}

El análisis hidráulico y de riesgos operativos para la profundización de las líneas de salmuera y condensado en el Humedal Arrieros ha permitido establecer las siguientes conclusiones técnicas:

\begin{table}[htbp]
    \centering
    \caption{Conclusiones técnicas derivadas del análisis hidráulico.}
    \label{tab:conclusiones}
    \setcounter{itemcount}{0}
    \begin{tabularx}{\textwidth}{@{}NX@{}}
        \toprule
        \multicolumn{1}{c}{\textbf{Item}} & \textbf{Conclusión} \\
        \midrule
        & El impacto hidráulico dinámico de la profundización es menor al \SI{2.0}{\percent} de variación en la pérdida de presión total del tramo. El incremento de presión dinámica es de tan solo \SI{0.33}{\kilo\pascal} (\SI{0.003}{bar}) para salmuera y \SI{0.39}{\kilo\pascal} (\SI{0.004}{bar}) para condensados, lo cual no genera afectación en la presión de succión de la estación Km8. \\
        & A condiciones operativas nominales de \SI{300000}{\kilogram\per\hour}, las velocidades de flujo en la tubería de \SI{12}{in} NPS (SDR 17) son de \SI{1.14}{\meter\per\second} (salmuera) y \SI{1.37}{\meter\per\second} (condensado). Ambas velocidades superan la velocidad crítica de autolimpieza de \SI{0.9}{\meter\per\second}, mitigando el riesgo de sedimentación de sólidos en suspensión durante la operación regular. \\
        & La configuración de sifón invertido de la PHD introduce un incremento localizado de presión estática en el fondo (cota \SI{2542.75}{m.s.n.m.}) que, sin válvulas de aislamiento, expone al tramo enterrado a la cota de origen de Sesquilé (\SI{2703.37}{m.s.n.m.}), resultando en una presión estática global máxima de \SI{18.90}{bar} para la línea de salmuera saturada y \SI{15.75}{bar} para la de condensado en condiciones de parada de planta. \\
        & La tubería original de HDPE PE100 SDR 17 (presión máxima admisible PN10 = \SI{10.0}{bar}) se encuentra severamente sobreesforzada bajo condiciones estáticas globales de parada sin aislamiento, superando su límite en un \SI{89.0}{\percent} (salmuera) y \SI{57.5}{\percent} (condensado). Cabe resaltar que la tubería superficial existente en el Km 16.5 ya presentaba presiones estáticas de \SI{17.07}{bar} en parada debido al desnivel global, excediendo también el límite nominal del SDR 17. \\
        & La geometría del perfil longitudinal crea dos crestas locales en los puntos de transición superficial (K0+000 y K0+450) propensas a la acumulación de bolsas de gas o aire libre, lo cual puede generar fenómenos de bloqueo por aire e inducir sobrepresiones mecánicas severas por golpe de ariete. \\
        & La auditoría de campo con el SCADA (Sección~\ref{sec:auditoria}) confirmó que la presión estática es independiente del caudal: es idéntica en la condición operativa actual y en el máximo futuro de \SI{300000}{\kilogram\per\hour}, por lo que el sobreesfuerzo del SDR 17 no se atenúa con el bajo flujo presente. Asimismo, evidenció válvulas de corte \texttt{ESDV-SAL-K17} y \texttt{ESDV-SAL-K8} ya instaladas, que habilitan la alternativa de aislamiento sin obra mayor, sujeto a verificar su posición respecto al cruce. \\
        \bottomrule
    \end{tabularx}
\end{table}

En síntesis, la alternativa de diseño por Perforación Horizontal Dirigida (PHD) es hidráulicamente viable para la succión de la estación Km8 en condiciones estacionarias nominales. Sin embargo, su viabilidad operativa global e integridad estructural a largo plazo depende de la mitigación de los riesgos mecánicos y operativos identificados en el fondo del sifón y en los puntos de transición superficial.
